\begin{textsample}{POS Dim 1 – to – Score 31.00 – to/to\_ef\_37.txt}  \label{ex:f1_pos_278}
Ciências da natureza O princípio da formação humana pressupõe a \textbf{integralidade} do \textbf{sujeito} e a sua \textbf{inter-relação} com o seu meio físico e social. \textbf{Enquanto} ser social, essa formação busca o seu despertamento ético para as questões básicas de convívio coletivo, e \textbf{enquanto} ser humano, a reaproximação com seu ambiente natural, o \textbf{que} \textbf{demanda} conhecimento para transformar e consciência para preservar.
Nesse processo de formação, as chamadas Ciências da Natureza objetivam tratar o desenvolvimento da Ciência como produto da construção histórica, social e cultural humana.
INTRODUÇÃO O letramento científico é o elemento principal de adensamento para as Ciências da Natureza, pois cria \textbf{um} elo entre as competências e habilidades.
O letramento científico propõe a construção de \textbf{um} conhecimento, \textbf{que} intervenha no mundo real e possa ser aplicado em diferentes situações cotidianas, \textbf{que} por sua vez, desenvolvam habilidades suficientes para a tomada de decisões com base nos procedimentos investigativos e no próprio desenvolvimento da ciência na história da \textbf{humanidade}.
Sob essa \textbf{ótica}, o letramento científico atua como articulador do conhecimento.
Nessa etapa de desenvolvimento da criança e do adolescente, deve -se assegurar aos estudantes o acesso aos conhecimentos científicos de modo \textbf{que} possam compreender, interpretar o mundo e agir sobre ele.
É importante enfatizar a necessidade permanente de flexibilização na organização curricular a partir da realidade de cada região do país, criando as condições necessárias para \textbf{que} a escola tenha autonomia em relação à construção de \textbf{um} currículo \textbf{que} atenda às necessidades dos estudantes.
Assim, a ciência se constrói a partir das necessidades humanas de compreender cada vez mais o ambiente em \textbf{que} se \textbf{vive} e de produzir conhecimento acerca daquilo \textbf{que} nos afeta de modo positivo ou negativo sejam as evoluções tecnológicas, sejam as doenças ou o nosso próprio modo de vida colocando -nos em situações complexas \textbf{que} \textbf{exigem} soluções inovadoras em relação ao uso dos recursos naturais. \textbf{Logo}, o ensino de Ciências nos anos iniciais do Ensino Fundamental é \textbf{um} importante momento de reflexão para a ação futura do cidadão em formação, \textbf{que} \textbf{necessita} apropriar -se do conhecimento já produzido pela \textbf{humanidade}, com vistas ao conhecimento \textbf{que} ainda pode ser produzido.
Por esse motivo, o processo de investigação, tão característico das Ciências da Natureza assume a \textbf{centralidade} da formação do cidadão do futuro e tem como premissa o letramento científico, elemento fundamental das finalidades pedagógicas do ensino de Ciências da Natureza. \textbf{Um} importante desafio com o qual nos deparamos é o de pensar o ensino de Ciências na perspectiva local de cada região e de cada comunidade.
O jovem estado do Tocantins, se por \textbf{um} \textbf{lado} ainda figura com parcos investimentos em desenvolvimento de ciência e tecnologia, tem por outro \textbf{uma} riqueza em recursos naturais inestimáveis nos ecossistemas mundiais.
As características do nosso ecossistema são motivos de pesquisas nacionais e internacionais, colocando -nos numa condição de lugar com muito a ser estudado e \textbf{conhecido}, permitindo \textbf{que} os próprios estudantes sejam investigadores de realidades ainda desconhecidas da ciência.
O Tocantins apresenta \textbf{uma} riqueza em biodiversidade e cultura inestimável para o Planeta.
Temos comunidades urbanas, rurais, indígenas, quilombolas, \textbf{ribeirinhas}, de quebradeiras de coco, pescadores, entre outras, e por isso apresentam \textbf{um} modo \textbf{bastante} particular de investigar seu cotidiano.
Essa riqueza nos traz também elementos para o trabalho com a ciência, a partir desses cotidianos diferenciados pela cultura e pelos modos de acumular conhecimentos, não \textbf{hierarquizando}, mas buscando nessas distintas trajetórias outros modos de observar, classificar, categorizar e analisar os fenômenos naturais \textbf{que} nos cercam.¹ ¹ Há para além do conhecimento acumulado pela ciência \textbf{um} vasto conhecimento \textbf{popular} para partir dele e desenvolver as habilidades e competências descritas na BNCC, valorizando temas \textbf{contemporâneos} \textbf{que} afetam a vida humana em escala local, regional e global.
Entre esses temas, destacam -se : direitos da criança e do adolescente ( Lei nº 8.069/1990¹⁶ ), educação para o trânsito ( Lei nº 9.503/1997¹⁷ ), educação ambiental ( Lei nº 9.795/1999, Parecer CNE/CP nº 14/2012 e Resolução CNE/CP nº 2/2012¹⁸ ), educação alimentar e nutricional ( Lei nº 11.947/2009¹⁹ ), processo de envelhecimento, respeito e valorização do idoso ( Lei nº 10.741/2003²⁰ ), educação em direitos humanos ( Decreto nº 7.037/2009, Parecer CNE/CP nº 8/2012 e Resolução CNE/CP nº 1/2012²¹ ), educação das relações étnico-raciais e ensino de história e cultura afro-brasileira, africana e indígena ( Leis nº 10.639/2003 e 11.645/2008, Parecer CNE/CP nº 3/2004 e Resolução CNE/CP nº 1/2004²² ), bem como saúde, vida familiar e social, educação para o consumo, educação financeira e fiscal, trabalho, ciência e tecnologia e diversidade cultural ( Parecer CNE/CEB nº 11/2010 e Resolução CNE/CEB nº 7/2010²³ ).
A Ciência, portanto, é capaz de produzir \textbf{equidade} num universo de relações díspares, promovendo, em cada canto do estado do Tocantins e do país, formas distintas de mobilizar conhecimentos para alcançar objetivos comuns.
Por esse motivo, as oito competências básicas para as Ciências da Natureza contemplam, de algum modo essa trajetória, como \textbf{veremos} \textbf{adiante} : 1.
Compreender as Ciências da Natureza como empreendimento humano, e o conhecimento científico como provisório, cultural e histórico. 2.
Compreender conceitos fundamentais e estruturas explicativas das Ciências da Natureza, bem como dominar processos, práticas e procedimentos da investigação científica, de modo a sentir segurança no debate de questões científicas, tecnológicas, socioambientais e do mundo do trabalho, continuar aprendendo e colaborar para a construção de \textbf{uma} sociedade \textbf{justa}, democrática e inclusiva. 3.
Analisar, compreender e explicar características, fenômenos e processos relativos ao mundo natural, social e tecnológico ( incluindo o digital ), como também as relações \textbf{que} se estabelecem entre eles, exercitando a curiosidade para fazer perguntas, buscar respostas e criar soluções ( inclusive tecnológicas ) com base nos conhecimentos das Ciências da Natureza. 4.
Avaliar aplicações e implicações políticas, socioambientais e culturais da ciência e de suas tecnologias para propor alternativas aos desafios do mundo \textbf{contemporâneo}, incluindo aqueles relativos ao mundo do trabalho. 5.
Construir argumentos com base em dados, evidências e informações confiáveis e negociar e defender ideias e pontos de vista \textbf{que} promovam a consciência socioambiental e o respeito a si próprio e ao outro, acolhendo e valorizando a diversidade de indivíduos e de grupos sociais, sem preconceitos de qualquer natureza. 6.
Utilizar diferentes linguagens e tecnologias digitais de informação e comunicação para se comunicar, acessar e disseminar informações, produzir conhecimentos e resolver problemas das Ciências da Natureza de forma crítica, significativa, reflexiva e ética. 7.
Conhecer, apreciar e cuidar de si, do seu corpo e bem-estar, compreendendo -se na diversidade humana, fazendo -se respeitar e respeitando o outro, recorrendo aos conhecimentos das Ciências da Natureza e às suas tecnologias. 8.
Agir pessoal e coletivamente com respeito, autonomia, responsabilidade, flexibilidade, resiliência e determinação, recorrendo aos conhecimentos das Ciências da Natureza para tomar decisões frente a questões científico-tecnológicas e socioambientais e a respeito da saúde individual e coletiva, com base em princípios éticos, democráticos, sustentáveis e solidários ( BRASIL, 2017 p. 322 ).
O desenvolvimento das competências específicas da área das Ciências da Natureza, somado aos procedimentos investigativos, refletem -se nas unidades temáticas Matéria e Energia, Vida e Evolução e Terra e Universo.
Quando falamos em procedimentos investigativos consideramos a intencionalidade do seu uso, como a construção de desafios, o questionamento acerca do seu cotidiano, o levantamento de dados, a análise e o tratamento desses dados, a comunicação dos resultados obtidos durante o processo de investigação e as intervenções necessárias para solucionar questões e modificar o meio em \textbf{que} se \textbf{vive} para contribuir com o desenvolvimento da sociedade.
Nesse sentido, para o desenvolvimento progressivo do estudante é importante valorizar a cooperação e o trabalho colaborativo, de modo a compartilhar as intervenções e a resolução de problemas, pilares da construção do fazer científico.
É importante destacar \textbf{que} essas investigações não estão necessariamente ligadas aos ambientes de laboratórios, mas aos mais diversos espaços \textbf{que} podem ser privilegiados na prática docente. \textbf{Ademais}, é \textbf{imprescindível} estimular o interesse e a curiosidade científica, oportunizando a definição de problemas do cotidiano, o levantamento de dados, o compartilhamento de resultados e ideias e a comunicação entre todos os envolvidos na atividade.
Esse conjunto de procedimentos \textbf{desencadeará} \textbf{um} ciclo contínuo de aprendizagem capaz de colocar o estudante diante de diferentes situações tornando -o capaz de propor soluções práticas no cotidiano, a fim de garantir a sua participação na resolução dos problemas.
Além disso, há de se destacar \textbf{que} para o desenvolvimento de \textbf{uma} educação científica no \textbf{atual} contexto, é \textbf{imprescindível} reconhecer os múltiplos papéis da tecnologia no desenvolvimento da sociedade humana.
É importante recuperar o debate sobre a ampliação das desigualdades e a degradação do ambiente em razão do mau uso do desenvolvimento tecnológico.
Para \textbf{que} esse diálogo seja possível nos diferentes contextos locais, a formação de professores deve abrir espaços para \textbf{que} o conhecimento acerca do desenvolvimento científico e tecnológico seja ampliado.
Desse modo, a fim de promover a \textbf{sistematização} do entendimento, neste documento, a área de Ciências da Natureza está organizada em três unidades temáticas, com base na Base Nacional Comum Curricular ( BNCC ), \textbf{que} progressivamente desenvolvem as habilidades e competências desejadas para esta etapa da educação básica, sendo elas : Matéria e Energia, Vida e Evolução e Terra e Universo.
Essas, articuladas ao letramento científico, trazem \textbf{uma} nova perspectiva do ensinar, pois possibilitam \textbf{que} a \textbf{interdisciplinaridade} e a integração com as demais ciências aconteçam a partir das situações cotidianas, com as características da investigação, muitas vezes necessárias para a compreensão do mundo \textbf{que} nos cerca.

% matched lemmas: ademais, adiante, atual, bastante, centralidade, conhecido, contemporâneo, demandar, desencadear, enquanto, equidade, exigir, hierarquizar, humanidade, imprescindível, integralidade, inter-relação, interdisciplinaridade, justo, lado, logo, necessitar, popular, que, ribeirinho, sistematização, sujeito, um, ver, viver, ótica
\end{textsample}
